\chapter{Literature Survey and Existing Systems}

\section{Review of Existing Systems}

Parking technology has evolved along four largely independent lines: mechanical
ticketing, vehicle-bound electronic tags, camera-based identification, and
IoT occupancy sensing. Each solves a real problem, and each leaves a specific gap
that shapes the design of \project{}.

\subsection{Manual ticketing and locally controlled barriers}

The conventional facility issues a ticket at entry, reads it at exit and charges at
a kiosk. The ticket is anonymous, so the system cannot associate a stay with a
person or a vehicle; revenue protection depends on the barrier and on staff. The
controller decides locally, which means a faulty loop detector, a lost ticket or a
power interruption becomes an operational incident rather than a query against a
record. This model is well understood and cheap per site, but it cannot support
account-based billing, history or remote operation without additional hardware.

\subsection{RFID tags and electronic toll collection}

Vehicle-bound RFID is the dominant model for tolling. India's National Electronic
Toll Collection programme issues a FASTag sticker that a fixed reader activates as
the vehicle approaches; the tag identifier is mapped to a prepaid or linked account
and the toll is deducted centrally \cite{npci2024}. The strengths are speed and
interoperability: the tag is passive, the reader is overhead, and the settlement is
centralised. The limitation for parking is structural. The tag identifies a
\emph{vehicle}, not a session, and the charge is a fixed toll for a crossing. A
parking fee depends on duration and on payment state, so an RFID-only design still
needs a separate session and billing system behind it. It also requires a reader at
every lane and a tag in every vehicle.

\subsection{Camera-based number-plate recognition}

ANPR systems photograph the plate, extract the characters and match them against a
record. They remove the tag but introduce a heavy dependency: illumination,
weather, plate formats and camera placement all affect accuracy, and each lane needs
a camera plus compute. Privacy is a further concern, because plate images are
personal data. The approach answers ``which vehicle is this?'' but not ``is this
driver allowed to leave, and has the fee been paid?'' --- those questions still need
an application layer.

\subsection{IoT occupancy sensing and smart-parking surveys}

A large body of work places a sensor at every bay --- ultrasonic, magnetometer,
infrared or camera --- and reports occupancy to a server. Recent surveys catalogue
these architectures, communication standards and business models
\cite{alam2023,biyik2021}. The value is guidance to a free space and utilisation
analytics. The cost is one device per bay, which dominates deployment economics for
a small facility, and the sensing problem is orthogonal to the access-control
problem: knowing that a bay is free does not tell a barrier whether to open.

\subsection{BLE beacons and phone-based parking}

BLE has been used mainly for proximity. Mackey et al.\ pair each parking spot with a
BLE beacon and use particle filtering on received signal strength to decide which
spot a phone has entered, with an automated payment scheme built on top
\cite{mackey2020}. The result demonstrates that commodity phones and low-cost
beacons are sufficient for reliable proximity decisions. It also concentrates the
intelligence on the phone and the cloud rather than on the barrier, and it assumes
one beacon per spot.

\subsection{Commercial mobile parking applications}

App-based parking payment (for example, phone applications that start and stop a
parking stay) is now common. These systems digitise the payment record but usually
leave enforcement and the barrier with the operator; the phone proves payment, and a
human or a camera checks it. They validate the central idea of \project{} --- that a
phone account is a usable parking credential --- while leaving the gate itself
unchanged.

\begin{table}[htbp]
  \centering
  \small
  \begin{tabularx}{\textwidth}{@{}lXXXX@{}}
    \toprule
    \textbf{Approach} & \textbf{Credential} & \textbf{Hardware per site} &
    \textbf{Billing basis} & \textbf{Main gap} \\
    \midrule
    Ticket + local barrier & Paper ticket & Barrier, dispenser, kiosk &
      Fixed tariff per stay & Anonymous; no account or history \\
    RFID / FASTag & Vehicle tag & Reader per lane, tags per vehicle &
      Fixed charge per crossing & Vehicle-bound; no duration model \\
    ANPR cameras & Number plate & Camera + compute per lane &
      Whatever the backend adds & Cost, accuracy, privacy \\
    Bay sensors & None & One sensor per bay & Time (computed centrally) &
      Does not control access; cost per bay \\
    BLE beacons (research) & Phone app & One beacon per spot &
      Usage time & Assumes per-spot beacons \\
    Mobile payment apps & Phone account & None at the gate &
      Time & Barrier and enforcement stay manual \\
    \textbf{\project{}} & \textbf{Phone account + JWT} & \textbf{One ESP32 +
      one servo} & \textbf{Per second, backend clock} &
      \textbf{Plaintext BLE; single gate (v1)} \\
    \bottomrule
  \end{tabularx}
  \caption{Comparison of existing approaches with the proposed system.}
  \label{tab:comparison}
\end{table}

\section{Gaps Identified and Proposed Solution}

Table~\ref{tab:comparison} makes the gaps concrete. Five matter for this project.

\begin{enumerate}
  \item \textbf{The barrier is trusted too much.} In ticketing, RFID and camera
        systems the local controller decides whether to open. \project{} moves the
        decision to the backend and reduces the controller to a courier with a
        device credential. A compromised phone can annoy the backend; it cannot open
        a gate.
  \item \textbf{Identity is attached to the wrong object.} Tags and plates identify
        vehicles, and tickets identify nothing. \project{} attaches identity to the
        account behind the phone, and then binds that account to a vehicle at
        session creation.
  \item \textbf{Duration is an afterthought.} Toll-style charging is per crossing,
        and bay sensors compute time but do not gate access. \project{} treats the
        duration as the primary quantity: the clock starts when the backend
        authorises entry and stops when payment succeeds.
  \item \textbf{Per-bay hardware does not scale down.} A small facility cannot
        justify a sensor at every bay. \project{} needs one ESP32 and one servo for
        the whole site and no per-bay device.
  \item \textbf{Verification is rarely reproducible.} Student and prototype
        projects often demonstrate a flow manually. \project{} ships a containerised
        verification workflow with a coverage gate and published reports, so the
        evidence can be regenerated by anyone with Docker.
\end{enumerate}

The proposed system therefore combines three ideas that exist separately in the
literature: BLE proximity discovery from research beacons \cite{mackey2020},
account-based settlement from tolling \cite{npci2024}, and server-side session
management from IoT parking platforms \cite{alam2023}. Its contribution is the
arrangement: a phone-carried session identifier, a single gate that re-validates
every action, and a lifecycle that makes duration and payment the only inputs to
the exit decision.


